\section{Fundamentação teórica}

\subsection{Arquiteturas de LLM relacionadas ao problema}

Os componentes mais diretamente relacionados a este trabalho são, de um lado, os LLMs generativos baseados em transformers \parencite{vaswani2017attention} e, de outro, os modelos de embeddings utilizados para recuperação semântica \parencite{nie2024textEmbeddingSurvey}. Enquanto o LLM atua como componente de síntese textual, os embeddings convertem documentos e consultas em vetores, permitindo o cálculo de similaridade semântica em bancos vetoriais \parencite{gao2023ragSurvey}.

Em sistemas RAG, esses componentes operam de maneira complementar: o modelo de embedding sustenta a recuperação, enquanto o modelo generativo formula a resposta a partir do contexto recuperado \parencite{gao2023ragSurvey}. Essa separação favorece uma arquitetura modular, a especialização de componentes e a evolução incremental da solução \parencite{gao2024modularRag}.

Segundo \textcite{gao2023ragSurvey}, sistemas RAG podem ser classificados em abordagens ingênuas, avançadas e modulares. Para o problema investigado, a arquitetura modular mostra-se adequada por permitir substituir, de forma independente, extratores documentais, estratégias de segmentação, modelos de embedding, bancos vetoriais, reranqueadores e modelos gerativos. Em contextos institucionais, essa característica é desejável porque infraestrutura, custos, segurança e políticas internas podem exigir alterações graduais da pilha tecnológica.

Essa modularidade, contudo, não deve ser interpretada como garantia automática de qualidade. A literatura apresenta o RAG como estratégia para reduzir problemas de desatualização, baixa verificabilidade e alucinação, mas não como solução suficiente para eliminá-los. O desempenho final continua dependente da qualidade documental, da segmentação, da recuperação, do desenho do prompt, da seleção de contexto e dos mecanismos de avaliação adotados \parencite{gao2023ragSurvey,gupta2024comprehensiveRag,saadFalconEtAl2024ares}. Essa leitura é importante em aplicações públicas, nas quais uma resposta plausível, mas mal fundamentada, pode comprometer confiança institucional.

\subsection{Técnicas relacionadas}

As principais técnicas adotadas neste projeto incluem extração documental, chunking, geração de embeddings semânticos, recuperação vetorial, reranqueamento, seleção contextual e prompting orientado a evidências \parencite{gao2023ragSurvey,maharjan2026chunkingReranking,wangEtAl2025segmentation}. Em conjunto, essas técnicas convertem o acervo em unidades recuperáveis, identificam trechos semanticamente relevantes e delimitam o contexto encaminhado ao modelo generativo. O prompting orientado a evidências complementa esse fluxo ao instruir o modelo a privilegiar o contexto recuperado, evitar inferências especulativas e apresentar referências verificáveis.

Dentre essas técnicas, o chunking ocupa posição central. Estudos recentes sobre segmentação documental em RAG indicam que escolhas inadequadas de divisão podem comprometer a recuperação semântica, seja por fragmentarem indevidamente o contexto, seja por gerarem blocos excessivamente extensos, com perda de precisão na busca \parencite{wangEtAl2025segmentation,bhatEtAl2025chunkSize}. No presente projeto, a segmentação documental exigiu sucessivos ajustes empíricos para equilibrar granularidade semântica, preservação de contexto e restrições operacionais do banco vetorial e da camada de ingestão. Assim, a experiência prática corrobora a literatura ao indicar que o chunking não constitui etapa meramente operacional, mas um dos principais determinantes do desempenho de um sistema RAG.

Surveys e frameworks recentes também indicam que, em ambientes reais, a qualidade de sistemas RAG depende de mecanismos de avaliação, rastreabilidade e comparação entre configurações, não apenas do pipeline básico de recuperação e geração \parencite{gao2023ragSurvey,saadFalconEtAl2024ares,esEtAl2023ragas}. Essa perspectiva orientou a criação de scripts próprios de preparação, importação, avaliação e documentação dos principais parâmetros experimentais.

No campo da avaliação, estratégias baseadas em LLM como juiz têm sido utilizadas para escalar a análise de respostas geradas, sobretudo quando combinadas com rubricas explícitas \parencite{liuEtAl2023geval,liuEtAl2025judgeRag}. Ainda assim, o julgamento automatizado pode variar conforme modelo avaliador, rubrica e abertura da tarefa; por isso, em aplicações RAG, tende a ser mais robusto quando tratado como instrumento auxiliar, combinado a evidências qualitativas, validação humana e análise das fontes recuperadas.

\subsection{Trabalhos e soluções relacionadas}

\textcite{gao2023ragSurvey} oferecem a base conceitual mais abrangente para compreender o RAG como resposta aos problemas de desatualização, baixa explicabilidade e geração não rastreável em LLMs. Seu trabalho orienta a decomposição do problema em etapas de recuperação, aumento de contexto e geração, fornecendo referencial teórico para o desenho da solução adotada neste estudo. No plano arquitetural, \textcite{gupta2024comprehensiveRag} consolidam o estado da arte em retrieval-augmented generation (RAG), discutindo como mecanismos de recuperação e geração se combinam para mitigar limitações de LLMs e destacando desafios de precisão, fidelidade e coordenação entre componentes.

Os trabalhos de \textcite{minaee2024llmSurvey} e \textcite{vaswani2017attention} situam a emergência dos grandes modelos de linguagem baseados em transformers e descrevem suas capacidades de compreensão e geração em ampla gama de tarefas, fornecendo o pano de fundo técnico para o uso de LLMs generativos neste projeto. \textcite{nie2024textEmbeddingSurvey} complementam esse quadro ao tratar os text embeddings como tecnologia fundacional para tarefas de busca semântica e recuperação de informação, enfatizando a relevância da integração entre LLMs e embeddings em aplicações contemporâneas.

A literatura recente sobre alucinação em LLMs evidencia que esses modelos podem produzir saídas fluentemente redigidas, porém factualmente incorretas ou não verificáveis, o que reforça, em contextos públicos sensíveis, a necessidade de soluções explicáveis e responsáveis, com mecanismos de fundamentação documental, governança e controle institucional \parencite{anhHoang2025hallucinations,gao2023ragSurvey,deAlmeidaSantos2025aiGovernance}.

Em termos de contexto institucional, \textcite{un2015responsiveGovernance}, \textcite{ipu2024aiParliaments}, \textcite{wfd2023guidelinesAiParliaments} e \textcite{geunis2023parliamentaryMonitoring} convergem ao apontar que parlamentos e administrações públicas operam sob pressões por transparência, responsividade, responsabilização e uso qualificado da informação. Nesse domínio, discursos e debates são insumos centrais para análise política e apoio técnico, mas os sítios parlamentares frequentemente funcionam mais como repositórios informacionais do que como instrumentos de exploração substantiva, deixando o acesso dependente de conhecimento prévio sobre estrutura e vocabulário \parencite{skubicFiser2024discourseReview,bandeiraBernardes2016information}.

No caso brasileiro, os Textos para Discussão da Consultoria Legislativa do Senado oferecem evidência empírica recente sobre a relevância e a complexidade do acervo de pronunciamentos. \textcite{blanco2025plenarioPalanqueEstudio} analisa discursos em plenário entre 2007 e 2024 e identifica mudanças de volume, extensão, tempo de fala e interatividade, associadas a transformações institucionais e tecnológicas. Em estudo complementar, \textcite{blanco2026letrasGraudas} examina leiturabilidade e traços morfossintáticos do mesmo domínio discursivo por meio de técnicas computacionais de análise textual. Esses trabalhos reforçam que os discursos do Senado constituem corpus relevante para métodos computacionais, embora se concentrem em análise descritiva e interpretativa, e não na construção de uma camada interativa de consulta, recuperação semântica e geração de respostas verificáveis.

No campo da governança pública de inteligência artificial, estudos e diretrizes institucionais enfatizam responsabilidade, transparência, monitoramento, gestão de riscos e alinhamento a finalidades públicas \parencite{deAlmeidaSantos2025aiGovernance,ipu2024aiParliaments,wfd2023guidelinesAiParliaments}. Para aplicações RAG em acervos legislativos, essa orientação reforça a importância de bases documentais controladas, metadados preservados, respostas verificáveis e avaliação contínua.

No domínio político, \textcite{khaliq2024ragar} demonstram que arquiteturas RAG podem apoiar tarefas de fact-checking político com raciocínio fundamentado em evidências externas, articulando recuperação rigorosa e geração explicável. Ainda que o objetivo deste trabalho não seja a verificação factual multimodal, a pesquisa desses autores é relevante por demonstrar a aderência do domínio político-parlamentar a abordagens baseadas em recuperação e geração com evidências.

Por fim, \textcite{martelloteViola2025organizacaoConhecimento} aproxima a discussão do contexto brasileiro ao tratar da inteligência artificial na organização do conhecimento legislativo. Sua contribuição reforça a pertinência de soluções explicáveis, documentalmente ancoradas e informacionalmente estruturadas em ambientes parlamentares, oferecendo base teórica específica para a aplicação proposta neste trabalho.

\subsection{Síntese da revisão e lacuna de pesquisa}

A revisão realizada permite identificar três conjuntos de literatura complementares. O primeiro, voltado a RAG, LLMs, embeddings, segmentação documental e avaliação automatizada, oferece a base técnica do trabalho. O segundo, dedicado a parlamentos, organização da informação legislativa e discurso parlamentar, delimita a relevância institucional do domínio. O terceiro, mais próximo do caso brasileiro, aplica métodos computacionais ao corpus de discursos do Senado, mas com foco predominante em análise descritiva, linguística ou discursiva \parencite{blanco2025plenarioPalanqueEstudio,blanco2026letrasGraudas}. A lacuna que permanece, portanto, não está na ausência de estudos sobre RAG nem na ausência de estudos sobre discursos parlamentares, mas na articulação aplicada entre esses dois campos.

Nesse sentido, embora a literatura sugira convergência quanto ao potencial do RAG para elevar a precisão factual e a verificabilidade de sistemas baseados em LLM, tal potencial depende de decisões metodológicas situadas, e não apenas da adoção da arquitetura RAG em abstrato. Parte significativa da literatura técnica permanece concentrada em discussões conceituais amplas, cenários corporativos ou aplicações em domínios distintos do acervo legislativo brasileiro. Ainda que pesquisas recentes já explorem computacionalmente discursos do Senado sob perspectivas quantitativas, discursivas e linguísticas, são menos frequentes estudos que articulem, em um mesmo desenho aplicado, corpus parlamentar nacional, preparação reprodutível de base documental, preservação de metadados legislativos, fluxo RAG operacional e avaliação empírica das respostas geradas.

Assim, observa-se lacuna aplicada relacionada à construção, documentação e avaliação de soluções RAG voltadas especificamente à consulta de discursos parlamentares no contexto do Senado Federal. Este trabalho busca contribuir para seu preenchimento ao propor, implementar e avaliar uma prova de conceito reprodutível de chatbot legislativo baseada em RAG, com foco em discursos da 56\textordfeminine{} Legislatura. A contribuição reside não apenas no protótipo, mas também na explicitação do pipeline, no enriquecimento dos chunks com metadados, na documentação dos artefatos de ingestão e avaliação, e na combinação de estudos de caso, rodadas automatizadas com rubrica e análise manual amostral.
